\begin{lemma}
\label{lemma-powers-field}
\begin{reference}
\cite[Lemma 3.1.6]{Alper-adequate}
\end{reference}
Let $k'/k$ be a field extension. The following are equivalent
\begin{enumerate}
\item for each $x \in k'$ there exists an $n > 0$ such that $x^n \in k$, and
\item $k' = k$ or $k$ and $k'$ have characteristic $p > 0$ and
either $k'/k$ is a purely inseparable extension or
$k$ and $k'$ are algebraic extensions of $\mathbf{F}_p$.
\end{enumerate}
\end{lemma}

\begin{proof}
Observe that each of the possibilities listed in (2) satisfies (1).
Thus we assume $k'/k$ satisfies (1) and we prove that we are in
one of the cases of (2). Discarding the case $k = k'$ we may assume
$k' \not = k$. It is clear that $k'/k$ is algebraic.
Hence we may assume that $k'/k$ is a nontrivial finite extension.
Let $k'/k'_{sep}/k$ be the separable subextension
found in Fields, Lemma \ref{fields-lemma-separable-first}.
We have to show that $k = k'_{sep}$ or that $k$ is an algebraic over
$\mathbf{F}_p$. Thus we may assume that $k'/k$
is a nontrivial finite separable extension and we have to show
$k$ is algebraic over $\mathbf{F}_p$.

\medskip\noindent
Pick $x \in k'$, $x \not \in k$. Pick $n, m > 0$ such that
$x^n \in k$ and $(x + 1)^m \in k$. Let $\overline{k}$ be an
algebraic closure of $k$. We can choose embeddings
$\sigma, \tau : k' \to \overline{k}$ with $\sigma(x) \not = \tau(x)$.
This follows from the discussion in
Fields, Section \ref{fields-section-separable-extensions}
(more precisely, after replacing $k'$ by the $k$-extension
generated by $x$ it follows from
Fields, Lemma \ref{fields-lemma-count-embeddings}).
Then we see that $\sigma(x) = \zeta \tau(x)$ for some
$n$th root of unity $\zeta$ in $\overline{k}$.
Similarly, we see that $\sigma(x + 1) = \zeta' \tau(x + 1)$
for some $m$th root of unity $\zeta' \in \overline{k}$.
Since $\sigma(x + 1) \not = \tau(x + 1)$ we see $\zeta' \not = 1$.
Then
$$
\zeta' (\tau(x) + 1) =
\zeta' \tau(x + 1) =
\sigma(x + 1) =
\sigma(x) + 1 =
\zeta \tau(x) + 1
$$
implies that
$$
\tau(x) (\zeta' - \zeta) = 1 - \zeta'
$$
hence $\zeta' \not = \zeta$ and
$$
\tau(x) = (1 - \zeta')/(\zeta' - \zeta)
$$
Hence every element of $k'$ which is not in $k$ is algebraic over the prime
subfield. Since $k'$ is generated over the prime subfield by the elements
of $k'$ which are not in $k$, we conclude that $k'$ (and hence $k$)
is algebraic over the prime subfield.

\medskip\noindent
Finally, if the characteristic of $k$ is $0$, the above leads to a
contradiction as follows (we encourage the reader to find their own proof).
For every rational number $y$ we similarly get a root of unity
$\zeta_y$ such that $\sigma(x + y) = \zeta_y\tau(x + y)$.
Then we find
$$
\zeta \tau(x) + y = \zeta_y(\tau(x) + y)
$$
and by our formula for $\tau(x)$ above we conclude
$\zeta_y \in \mathbf{Q}(\zeta, \zeta')$. Since the number field
$\mathbf{Q}(\zeta, \zeta')$ contains only a finite number of roots of
unity we find two distinct rational numbers $y, y'$ with
$\zeta_y = \zeta_{y'}$. Then we conclude that
$$
y - y' =
\sigma(x + y) - \sigma(x + y') =
\zeta_y(\tau(x + y)) - \zeta_{y'}\tau(x + y') = \zeta_y(y - y')
$$
which implies $\zeta_y = 1$ a contradiction.
\end{proof}